%!TEX root = ../../thesis.tex

\subsection{Reward and Episode Structure}\label{sec:input-reward}

% TODO: Reward first, episode second, then the interaction between them.
%
% Reward: the living reward for surviving a step, the terminal coverage burst, and
% the novelty bonus --- what each one is for, and the scale each is given. Show the
% arithmetic that ties the reward magnitudes to the categorical value support
% (the +-1 step-reward gap must span more than one bin or the value and value-prefix
% heads go action-blind); this is the worked example of a hyperparameter *derived*
% from problem structure rather than tuned, and \cref{sec:methodology-hyperparameters}
% refers back to it.
%
% Episode: an episode ends on early death at a wrong character or on the fixed input
% length, whichever comes first. Reset is the snapshot restore of
% \cref{sec:platform-snapshots}. Give the horizon actually used and say why it is
% bounded tightly (long or absorbing horizons feed the dead-branch failure mode of
% \cref{sec:failure-dead-branch}).
%
% Interaction: the reward under this formulation is hand-shaped, and that is exactly
% what \cref{sec:discussion-threats} identifies as the strongest objection to the
% result. Say so here rather than letting the reader discover it, and draw the
% contrast with the corpus formulation (\cref{sec:corpus-reward}), where the reward is
% raw coverage and was not designed around the agent.
%
% Reward hacking belongs here as a designed-against failure, not an anecdote: the
% teardown burst made dying profitable, which is why the novelty bonus is gated on
% survival. Cross-reference \cref{sec:failure-reward-hacking}.
